Optical lattice and single-ion clocks now reach fractional systematic
uncertainties at or below \selfEvalBound~\citep{mcgrew2018geodesy,
aeppli2024sr, lu2025ntsc, marshall2025alplus, jia2026sr, arnold2026lu,
zhang2026ca, zhu2026yb}, two orders of
magnitude better than current primary frequency standards~\citep{heavner2014nistf2,
weyers2018ptbcsf}. Comparing independent clocks in different regions
and laboratories verifies these self-assessed precision budgets more
objectively: the environmental conditions and geodetic assumptions
entering a self-evaluation are typically tied to a single location, and
identical methods within one laboratory may share common systematic
biases. Linking clocks of different types and locations into a
collaborative network is therefore a key pathway toward establishing
optical clocks as a verifiable metrological
infrastructure~\citep{komar2014quantumnetwork, riehle2017networks}. The
idea has moved from proposal to practice: fiber links with noise levels
below those of the clocks themselves now span thousands of
kilometers~\citep{schioppo2022cavities, chen2026dissemination}, and
multi-clock campaigns have measured tens of ratios among clocks in six
countries~\citep{lindvall2025coordinated} and across a continental
network~\citep{pizzocaro2026european}. Satellite-based advanced links
with comparable performance~\citep{shen2021satellite, shen2022freespace,
caldwell2023geo, derevianko2022spaceclock} are also being developed and
will interconnect with fiber networks to form a future global optical
clock network.

The most immediate purpose of such a network is the redefinition of the
SI second. One of the mandatory requirements (I.2.b) in the roadmap for
this redefinition~\citep{dimarcq2024roadmap} is that at least five
non-unit frequency ratios between standards of different clock
transitions be measured, each ratio at least twice by different
institutes, in agreement with an overall uncertainty of the comparison
$\Delta\nu/\nu < \refThreshold$---a requirement reiterated in the 2025
CCTF progress report~\citep{cctf2025strategy, cctf2025taskforce}.
However, comparisons across independent institutions and different
locations have reached at best a precision of
\refEuBestU~\citep{pizzocaro2026european}. So far, no frequency-ratio
measurements between clocks in different cities have satisfied this
uncertainty requirement~\citep{fortier2026review}.

Closing that gap matters beyond the definition itself. Chronometrically
leveled clocks map the gravitational potential of the Earth at
centimeter resolution~\citep{grotti2018geodesy, mcgrew2018geodesy,
takamoto2020redshift}, and clock networks have been proposed and used as
sensors for dark matter~\citep{derevianko2014topological,
wcislo2018global, beloy2021bacon}. All of these applications assume
that remote clocks agree at a level that, until now, has been
demonstrated only locally, in three experiments: a Sr/Yb/Al$^+$
triple-clock ratio measurement with an uncertainty of
\refSameCampusU~by NIST--JILA in Boulder (US)~\citep{aeppli2026nist},
and two transportable-clock comparisons at the $10^{-18}$
level~\citep{takamoto2020redshift,
vishwakarma2026interspeciesclockcomparison5}.

Here we report an intercity frequency-ratio measurement
between an IAPMST ytterbium lattice clock in Wuhan and a USTC strontium
lattice clock in Shanghai, about $\siteKm$\,km apart, connected by a
$\fibreKm$\,km fiber link. We determine the Yb/Sr ratio with a total
uncertainty of \uTotalE$\times10^{-18}$; gravitational redshift
effects have been corrected, whereas measurement error and the
time-varying gravitational redshift term associated with gravitational
tides have been retained as uncertainty components. The result
establishes the two sites as a prototype node pair for a distributed
optical-clock network, and demonstrates, at intercity scale, the
class of ratio measurement that the redefinition process depends on.
